% !TEX program = xelatex
% ---------------------------------------------------------------------------
% 北京语言大学本科生毕业论文
% 组成顺序：封面→题名页→原创声明→中英文摘要→目录→正文→结论→参考文献→附录→致谢
% 参考文献另起一页；致谢在最后。
% ---------------------------------------------------------------------------
\documentclass{base}
\usepackage{amssymb}
\usepackage{amsmath}
\usepackage{graphicx}
\DeclareGraphicsExtensions{.png,.pdf,.jpg,.jpeg}
\usepackage{float}
\usepackage{tikz}
\usetikzlibrary{arrows.meta,positioning,shapes.geometric,calc,fit,decorations.pathreplacing,matrix}
\usepackage{gbt7714}
\usepackage{silence}
\WarningFilter{caption}{Unused}

% ====================================================
% 论文基本信息
% ====================================================
\title{情感交互Live2D数字人系统设计}
\entitle{Design and Implementation of a Live2D Digital Avatar System with Affective Interaction}
\author{钟文源}
\enauthor{Wenyuan Zhong}
\stunum{202211690474}
\grade{2022级}
\college{信息科学学院}
\encollege{School of Information Science}
\major{人工智能}
\instructor{汪柳萍}
\eninstructor{Liuping Wang}
\degreetitle{Bachelor of~Engineering}
\date{二〇二六}{六}{2026}

\begin{document}

\maketitle
\tableofcontents

% ====================================================
% 中文摘要
% ====================================================
\begin{cnabstract}
随着高校心理健康教育需求的提升与虚拟数字人技术的普及，传统纯文本对话系统缺乏情感反馈、Live2D演示工程缺乏完整后端支撑的问题日益突出。本文面向高校课程演示与学生低压力情感交互练习场景，设计并实现了一套基于浏览器的、须登录使用的Live2D数字人Web系统，核心解决``形象驱动、语音播报、分层记忆持久化与时间驱动主动关怀''四大需求，以工程可复现仓库形式交付，明确不用于医疗或心理咨询，不替代专业人员服务。

系统采用B/S架构，前端基于Vite与Live2D Cubism Web构建，负责角色渲染、交互操作及私有资源预签名加载；语音经语音识别网关对接云端DashScope Fun-ASR完成识别后，与文本输入汇合，统一经对话WebSocket接口上行至服务端。服务端以FastAPI为核心，采用Ollama双模型分工模式：主对话模型（Qwen2.5系列）负责流式自然语言生成，动作决策模型输出结构化表情与动作字段，通过资源目录索引约束确保可执行性；语音合成默认采用云端小米MiMo服务，保留本地GPT-SoVITS切换路径。数据持久化层由关系型数据库、Redis分层缓存及对象存储构成，服务端通过统一REST接口集提供标准化响应。

记忆系统实现瞬时、短期、长期三级分层架构：瞬时记忆以低时延缓存保留最近若干轮对话原文并直接注入对话上下文；短期记忆对历史轮次进行规则压缩后持久化并拼入系统提示；长期记忆由后台任务周期性从历史会话中生成概要文本并缓存以备快速读取。此外，用户画像在每隔若干轮对话或断线时由LLM自动合并近期对话记忆刷新，注入主对话语境。主动关怀系统实现了从``对话内容抽取—关怀信息入库—后台定时调度—个性化话术生成—即时推送''的完整闭环：每轮对话结束后系统异步抽取结构化关怀信息并与关联对话上下文绑定；常驻后台调度器（默认300秒间隔）扫描到期记录，结合角色人设与用户画像生成个性化话术，经WebSocket推送文本及可选语音，离线用户于重连后自动补发。

在开发环境下，系统主链路已联调贯通，覆盖了登录、文本与语音对话、Live2D模型切换、历史会话拉取、模型资源入库、分层记忆读写及主动关怀推送等核心功能。记忆分层设计借鉴阿特金森--谢弗林多层存储思路与MemGPT核心思想，由服务端固定路径统一编排读写时机；主动关怀机制将时间感知与长期记忆相结合，使数字人系统初步具备超越单轮被动应答的主动情感关怀能力。

    \cnkeywords{情感交互；虚拟数字人；Live2D；FastAPI；分层记忆；主动关怀；WebSocket}
\end{cnabstract}

% ====================================================
% 英文摘要
% ====================================================
\begin{enabstract}
With the increasing demand for mental health education in universities and the popularization of virtual digital human technology, traditional pure-text dialogue systems increasingly fall short in providing emotional feedback, while Live2D demonstration projects commonly lack complete back-end support. Targeting the scenarios of college course demonstrations and students' low-pressure emotional interaction practice, this thesis designs and implements a browser-based, login-required Live2D digital avatar Web system. The system addresses four core requirements: image-driven interaction, voice broadcast, persistent hierarchical memory, and time-driven proactive care. It is delivered as a reproducible engineering repository and is explicitly not intended for medical or psychological counseling, nor does it replace professional services.

The system adopts a B/S architecture. The front-end is built on Vite and Live2D Cubism Web, responsible for character rendering, interactive operations, and presigned loading of private resources. Voice input is recognized via the \texttt{/ws/asr} interface connected to the cloud DashScope Fun-ASR service, then merged with text input and sent to the server through the dialogue WebSocket interface (\texttt{/ws/chat}). The server core is FastAPI, employing a dual-model Ollama collaboration (Qwen2.5 series): the main dialogue model handles streaming natural language generation, while the action decision model outputs expression and motion fields constrained by a resource directory index to ensure executability. Speech synthesis defaults to the cloud Xiaomi MiMo service, with a local GPT-SoVITS alternative retained. The data persistence layer comprises relational database tables, Redis layered cache, and object storage, with unified REST interfaces providing standardized responses.

The memory system implements a three-tier hierarchical architecture: instant memory preserves the most recent $K$ conversation turns in low-latency cache and injects them directly into the dialogue context; short-term memory applies rule-based compression to evicted turns and appends them to the system prompt; long-term memory is periodically synthesized from historical session records by a background task and cached for rapid retrieval. Additionally, the user profile is automatically refreshed by an LLM integrating recent dialogue memory every $N$ turns or upon disconnection. The proactive care system implements the complete pipeline from in-dialogue information extraction, care-trigger persistence, background scheduling, and personalized speech generation to real-time delivery: after each conversation concludes, the system asynchronously extracts structured care information bound to the associated dialogue context; a persistent background scheduler (default 300-second interval) claims due records and generates personalized care messages combining persona, user profile, and contextual conversation, then pushes text and optional MiMo speech via WebSocket, with offline users receiving pending messages upon reconnection.

In the development environment, the system's main pipeline is fully integrated and covers core functionality including login, text and voice dialogue, Live2D model switching, session history retrieval, model resource import, hierarchical memory read/write, and proactive care delivery. The memory hierarchy draws on the Atkinson-Shiffrin multi-store model and MemGPT's core ideas, with the server orchestrating read/write timing through fixed code paths. The proactive care mechanism combines temporal awareness with long-term memory, enabling the digital avatar system to transcend passive single-turn responses and achieve active emotional care capabilities.

    \enkeywords{affective interaction; digital human; Live2D; FastAPI; layered memory; proactive care; WebSocket}
\end{enabstract}

% ====================================================
% 正文开始
% ====================================================
\clearpage
\mainbodytrue
\pagenumbering{arabic}
\pagestyle{mainstyle}

% ====================================================
% 第一章 引言
% ====================================================
\chapter{引言}

\section{研究背景与意义}

\subsection{研究背景}

我国高校心理健康工作受到政策持续推动，教育部等部门的专项行动计划强调完善心理健康教育与咨询服务体系\cite{MOE2023}；国民心理健康发展报告从群体调查角度呈现青年阶段情绪与压力管理需求\cite{NMHReport2023}。世界卫生组织对精神卫生服务的倡导提示，可扩展、可触达的支持形式有助于弥补专业资源在地区与时段分配上的不均衡\cite{WHO2022}。需要强调的是，\textbf{数字人陪伴系统不能替代专业心理干预}，但在``日常情绪表达、低压力对话、习惯化自我梳理''等边界内，具有可验证的工程价值\cite{Picard1997,Bickmore2005}。

与上述需求相对照，\textbf{通用大模型聊天页}往往停留在``纯文本+单轮/弱连续''的交互形态；以Character.AI、Replika等为代表的\textbf{商业情感陪伴应用}体验完整但可复现性弱、难与课程设计深度结合\cite{CharacterAIAbout,ReplikaOfficial}；\textbf{大量Live2D示例工程}又常止步于资源展示，缺少与LLM、TTS、账户体系、分层记忆、主动关怀一体的可扩展后端\cite{live2d}。因此，有必要在Web端构建一条\textbf{可运行、可迭代、可追溯验证}的完整链路，使``看得见形象、听得到声音、存得下记忆、关怀得到用户''同屏实现。

\subsection{研究意义与差异化定位}

本文意义不在于提出全新生成模型，而在于\textbf{面向可控教学/研究场景}，给出可交付的系统化集成。其一，浏览器侧以\textbf{对话WebSocket接口}（\texttt{/ws/chat}）合并下行文本与朗读WAV，语音输入另经\texttt{/ws/asr}上行，职责分离且连接数可控。其二，Ollama\textbf{双模型分工}将``说话内容''与``表情/动作资源ID''解耦，便于分别调参与排错。其三，三级分层记忆在\textbf{Redis+MySQL}上完整落地，实现会话数据的有效沉淀与复用。其四，主动关怀系统完成从对话抽取到WebSocket推送的\textbf{全链路闭环}，使系统具备超越被动应答的时间感知与主动关怀能力。

\subsection{研究目的、范围与实验条件}

\textbf{目的与范围：}在浏览器内交付一套可运行的Live2D数字人Web系统，覆盖登录与会话、文本与语音对话、流式文本与音频下发、模型与资源管理、分层记忆持久化、主动关怀推送等核心能力；明确不覆盖医疗诊断、心理咨询资质替代及原生客户端开发。

\textbf{理论依据与实验条件：}记忆分层在工程上与经典多层存储模型及MemGPT类体系作类比说明\cite{Atkinson1968,MemGPT2023}。实验与联调环境为常见x86开发机，依赖现代浏览器、FastAPI宿主、Ollama推理服务、可配置的TTS与ASR密钥，以及MySQL、Redis、MinIO等按文档部署的存储服务；与性能相关的表述仅基于开发环境下的非正式观测，不作严格控制实验宣称。

\subsection{本文章节安排}

图~\ref{fig:layer_arch}给出由浏览器、应用服务、模型与媒体至持久化的\textbf{四层逻辑架构}。后续章节安排：第二章归纳关键技术与选型对比；第三章为需求分析与总体设计；第四章按前端实现、WebSocket通信、服务端编排、数据持久化与API四块阐述核心模块；\textbf{第五章}重点论述分层记忆系统；\textbf{第六章}重点论述主动关怀系统；第七章给出结论、创新点与未来工作。

\begin{figure}[htbp]
    \centering
    \includegraphics[width=\textwidth]{fig_layer_arch.png}
    \caption{系统四层逻辑架构示意图（UML组件图风格）}
    \label{fig:layer_arch}
\end{figure}

\section{国内外研究现状}

\subsection{陪伴型对话系统}

国外关于陪伴型对话系统的研究与产品演进较早，形成了从关系型智能体理论到大规模应用平台的连续谱系。早期研究强调长期关系建立与信任维持机制\cite{Bickmore2005}，随后在心理健康支持方向出现了以Woebot为代表的可量化评估工作\cite{Fitzpatrick2017}。Replika与Character.AI分别代表``长期情感陪伴''与``角色化对话社区''两类典型路径，二者都强调持续互动、个体化角色与多轮上下文体验\cite{ReplikaOfficial,CharacterAIAbout}。小冰体系将同理心计算与长期会话指标工程化，提供了陪伴型系统从实验到规模化运营的标志性案例\cite{ShumXiaoice2020}。相关研究进一步指出，陪伴型场景相较通用问答更依赖人格一致性、记忆连续性与安全边界设计\cite{Roller2021,Zhao2023Survey}。

国内工业界在``数字人+大模型+行业场景''融合方面推进较快。阿里云虚拟数字人平台、商汤如影与网易伏羲分别在形象定制、数字人直播与虚拟智能体领域形成产品化方案\cite{AliyunAvatarProduct,SenseAvatarProduct,NeteaseFuxiAgent}。学术侧以情感对话生成与共情策略建模\cite{Li2021EmotionDialogueCN,Chen2022EmpatheticDialogueCN}、多轮记忆增强\cite{Sun2022DialogueMemoryCN}、用户画像驱动的个性化服务\cite{He2023UserProfileCN}为主线。总体而言，学术工作在算法模型与应用讨论上积累较快，但与可复现工程系统之间仍有距离，尤其在``前端表现层—后端推理层—记忆治理层—主动关怀层''一体化验证方面仍需加强。

\subsection{LLM 驱动 Agent 与记忆机制}

大模型推动对话系统从``回复生成器''向``可规划、可工具调用、可长期记忆''的Agent形态迁移。ReAct（推理与行动协同框架，Reasoning and Acting）将推理与行动耦合到统一轨迹\cite{ReAct2023}；Generative Agents给出``观察—记忆—反思—计划''架构，验证了多智能体社会行为模拟的可行性\cite{GenerativeAgents2023}。MemGPT聚焦于``上下文长度有限、其余放进库里按需取用''这一核心问题，为工程层记忆分层提供了重要参照\cite{MemGPT2023}。检索增强生成（RAG）进一步扩展了知识调用边界\cite{Lewis2020}。在框架层，OpenCLaW、Harness等工具链进一步推动多Agent协作的评测标准化与工程化落地\cite{OpenCLaW2024,Harness2024}。

\subsection{现状小结与本文切入点}

综合文献与产业实践，可归纳出三点：一是陪伴型系统已从``可聊天''转向``可长期关系维护''；二是数字人生态从单点渲染能力转向角色运营与多模态内容生产；三是LLM Agent正在重塑系统架构，但工程稳定性与记忆治理仍是主线问题。本文的切入点是将上述几条线在Web工程中落地为可运行原型：以Live2D承担角色呈现，以双模型与WebSocket流式链路承担对话与朗读编排，以结构化记忆三层架构承担会话沉淀，以时间驱动调度器承担主动关怀闭环。

\section{研究内容与创新点}

本文以CubismDemo工程为基础，完成以下实现任务：一是\textbf{前端交互模块}，实现Live2D渲染接入、双栏交互界面、流式字幕展示、模型包与背景切换、语音输入与识别回填；二是\textbf{服务端会话编排模块}，基于FastAPI构建WebSocket会话服务，接入Ollama双模型，朗读在云端MiMo与本地GPT-SoVITS间可配置；三是\textbf{分层记忆模块}，设计并实现瞬时/短期/长期三级记忆架构，含Redis读写策略、规则压缩算法、周期概要后台任务与用户画像自动刷新；四是\textbf{主动关怀模块}，实现对话抽取、触发器调度、LLM话术生成、WebSocket推送的完整闭环；五是\textbf{数据与API模块}，设计MySQL八表并实现58条REST路由，统一\texttt{code/message/data}响应。

% ====================================================
% 第二章 关键技术基础与选型
% ====================================================
\chapter{关键技术基础与选型}

\section{Live2D Cubism 角色呈现}

Live2D通过二维网格变形与参数曲线驱动表情/姿态，配合\texttt{physics}、\texttt{motion}、\texttt{expression}等资源形成可渲染闭环\cite{live2d}。本项目将\textbf{expression/motion标识}限定于当前模型包目录索引（\texttt{catalog}），服务端校验后再驱动前端，避免``模型幻觉驱动不可执行资源''。相较VRM与三维引擎结合方案（如Three.js），Live2D以较低算力与成熟SDK换取可用的视觉反馈密度，契合``对话为主、形象为辅''的定位；相较纯SVG/CSS虚拟形象方案，则具备商用级口型与动作资产生态。

\section{WebSocket 全双工通信}

WebSocket在单一TCP连接上提供全双工帧传输，适合交错下发JSON事件与二进制音频载荷\cite{rfc6455,mdn-ws}。HTTP长轮询头部开销大、实时性差；SSE适合服务端单向文本流，但难以在同一语义通道内无损混合二进制音频帧。本项目采用WebSocket；\textbf{当前前端采用单连接合并流}：在\texttt{/ws/chat}上先发JSON文本帧，再在启用TTS时下发\texttt{chunk\_audio}元数据并紧随二进制WAV，主动关怀帧\texttt{type:remind\_trigger}亦经同一连接投递。

\section{Ollama 双模型分工}

本项目以Qwen2.5系列模型（通义千问2.5，Alibaba Qwen Team）作为主对话模型与动作决策模型的后端推理底座，借助Ollama\cite{ollama}本地托管服务统一调度。聊天模型以流式方式生成自然语言；动作模型以非流式调用输出结构化JSON，字段含\texttt{expression}与\texttt{motion}，由服务端解析、校验与回退。Qwen2.5在中文理解、指令遵循与格式化输出方面的表现，使其适配双模型分工场景：主对话侧要求上下文长、情感自然，动作侧要求格式严格且延迟可控。相较单模型混合生成，分工设计可将调参空间解耦，格式稳定性更高，调试日志更清晰；代价是一次交互需两次推理调用，总时延略有上升，本项目以工程可控性换取稳定性。

\section{语音合成：云端 MiMo 与本地 GPT-SoVITS}

无论后端指向本地引擎还是云端API，合成侧均先将回复文本按句末标点攒批切段，并行请求TTS，回收WAV后统一封装，再按音频序号保序经WebSocket下发。\textbf{本地合成}（GPT-SoVITS\cite{gptsovits}）将声学与算力留在自有机器，数据路径可控；代价是显存与算力与对话侧模型竞争，消费级单卡易出现排队与稳定性问题。\textbf{云端合成}（小米MiMo\cite{mimoV25Docs}）把声学推理外移，本机GPU主要留给对话与动作模型，弱显卡环境下端到端体验更稳定；代价是网络时延与供应商绑定。\textbf{工程结论}：TTS在代码层抽象为可插拔实现，通过配置在两条路径间切换，调用方与前端协议保持不变；\textbf{当前仓库默认以MiMo为主线}，GPT-SoVITS分支保留于代码树，但示例配置未随主线同步维护。

\section{数据三元结构}

MySQL负责结构化实体与审计友好CRUD\cite{mysql80}；Redis承担\textbf{毫秒级读写窗口}，覆盖瞬时与短期列表及长期合并字符串缓存\cite{redis}；MinIO卸载大文件对象（Live2D模型资源包、TTS参考音频等）\cite{minio}，通过预签名URL实现浏览器侧私有安全加载。三者分工明确：按需组合服务不同业务场景，无需在性能与可审计性之间单一取舍。

% ====================================================
% 第三章 系统需求分析与总体设计
% ====================================================
\chapter{系统需求分析与总体设计}

\section{需求分析}

\subsection{需求来源与边界说明}

本系统需求来自三类约束：用户侧需要低门槛、连续的文本、语音与Live2D形象同时可用的交互；开发侧需要可演进、可替换、可排障的工程结构；文档与实现之间需要边界清晰、验收口径可追溯。参与者划分为：\textbf{用户}（完成注册或登录）、\textbf{已登录用户}（对话、模型切换、接收关怀、查看历史）、\textbf{资源管理用户}（上传资源与配置人设）、\textbf{后台调度器}（定时扫描触发器并推送关怀）。

\subsection{UML 用例分析}

图~\ref{fig:usecase_overview}给出系统UML用例图。本系统不提供访客可用的业务入口，所有核心功能均在认证通过后方可使用。语音对话经识别后复用文本发送链路（\texttt{<<include>>}关系）；主动关怀推送由后台调度器通过\texttt{/ws/chat}投递，\textbf{已完整实现}。

\begin{figure}[H]
    \centering
    \includegraphics[width=0.96\textwidth]{fig_usecase.png}
    \caption{系统UML用例图（含系统边界；无访客模式；主动关怀闭环已完整实现）}
    \label{fig:usecase_overview}
\end{figure}

\subsection{功能需求明细}

表~\ref{tab:functional_requirements}给出功能需求明细，编号FR-xx在后文可作需求追溯索引。

\begin{table}[H]
    \centering
    \caption{系统功能需求明细表}
    \label{tab:functional_requirements}
    \small
    \setlength{\tabcolsep}{3.5pt}
    \begin{tabular}{|p{1.1cm}|p{2.2cm}|p{3.2cm}|p{2.3cm}|p{3.3cm}|}
        \hline
        \textbf{编号} & \textbf{功能名称} & \textbf{需求说明} & \textbf{输入/输出} & \textbf{验收要点} \\
        \hline
        FR-01 & 登录与会话初始化 & 用户完成身份进入，系统生成并持久化会话相关参数。 & 账号→会话状态 & 未登录无法进入主页面 \\
        \hline
        FR-02 & 文本对话 & 用户输入文本，服务端返回流式文本片段并逐段展示。 & message→chunk/done & 可见流式输出 \\
        \hline
        FR-03 & 语音对话 & 语音经ASR转为文本，复用文本对话主链路。 & 音频→识别文本+对话结果 & 结果与文本入口一致 \\
        \hline
        FR-04 & 角色动作与表情 & 依据结构化字段驱动Live2D表情/动作。 & expression/motion→动作 & 字段非法时有回退策略 \\
        \hline
        FR-05 & TTS语音播报 & 对话文本切段后触发TTS，经WebSocket分段下发。 & 文本段→chunk\_audio+WAV & 音频按序播放 \\
        \hline
        FR-06 & 模型与背景切换 & 用户可切换角色模型与场景背景。 & package→新模型状态 & 切换后无残留状态 \\
        \hline
        FR-07 & WebSocket异常恢复 & 网络中断时客户端自动重连。 & 断连→重连可用连接 & 延迟重连可触发 \\
        \hline
        FR-08 & 会话历史沉淀 & 每轮对话写入持久化存储，支持分页查询。 & 消息→会话记录 & 历史可按用户/分页读取 \\
        \hline
        FR-09 & 分层记忆读写 & 瞬时/短期/长期三级记忆在Redis+MySQL完整落地。 & 对话→记忆更新→注入 & 三层均可追溯验证 \\
        \hline
        FR-10 & 主动关怀推送 & 基于触发器规则在到期时生成个性化话术并推送。 & 触发规则→WebSocket帧 & 已实现全链路闭环 \\
        \hline
        FR-11 & Live2D资源上传与索引 & 上传模型ZIP，完成解析、对象存储与数据库索引。 & ZIP→package索引 & 可列出包信息并删除 \\
        \hline
        FR-12 & REST统一接口 & 提供61条\texttt{/api/*}路由，统一code/message/data响应。 & HTTP请求→结构化响应 & 异常场景返回可读信息 \\
        \hline
    \end{tabular}
\end{table}

\section{总体架构与主对话时序}

系统采用B/S结构，图~\ref{fig:seq_main_dialogue}给出主对话端到端时序，涵盖键盘与语音双入口汇合、Ollama双模型调用、TTS切段保序下发与记忆轮末写回等关键步骤。

\begin{figure}[htbp]
    \centering
    \resizebox{\textwidth}{!}{
    \begin{tikzpicture}[
        font=\small,
        participant/.style={rectangle, draw=black!70, thick, minimum width=2.1cm, minimum height=0.6cm, align=center, fill=white},
        lifeline/.style={dashed, thick, gray!50},
        msg/.style={-{Latex[length=2mm]}, thick},
        retmsg/.style={-{Latex[length=2mm]}, thick, dashed},
        note/.style={rectangle, draw=black!35, fill=yellow!10, rounded corners=2pt, font=\scriptsize, inner sep=2.5pt, align=left},
    ]
        \node[participant] (U)   at (0,   0) {用户};
        \node[participant] (BR)  at (3.2, 0) {浏览器\\[-2pt]{\scriptsize ws.js}};
        \node[participant] (ASR) at (6.2, 0) {ASR网关\\[-2pt]{\scriptsize /ws/asr}};
        \node[participant] (DS)  at (9.0, 0) {DashScope\\[-2pt]{\scriptsize Fun-ASR}};
        \node[participant] (WS)  at (12.0,0) {对话编排\\[-2pt]{\scriptsize /ws/chat}};
        \node[participant] (OL)  at (15.3,0) {Ollama\\[-2pt]{\scriptsize 双模型}};
        \node[participant] (TTS) at (18.2,0) {TTS服务\\[-2pt]{\scriptsize MiMo}};
        \node[participant] (DB)  at (21.0,0) {Redis/MySQL};

        \foreach \x in {0,3.2,6.2,9.0,12.0,15.3,18.2,21.0}
            \draw[lifeline] (\x,-0.38) -- (\x,-16.5);

        \node[note] at (1.6,-1.0) {\textit{路径A：语音}};
        \draw[msg] (0,-1.4) -- node[above,font=\scriptsize]{语音输入} (6.2,-1.4);
        \draw[msg] (6.2,-1.8) -- node[above,font=\scriptsize]{PCM二进制帧} (9.0,-1.8);
        \draw[retmsg] (9.0,-2.2) -- node[above,font=\scriptsize]{识别中间/终稿事件} (6.2,-2.2);
        \draw[retmsg] (6.2,-2.6) -- node[above,font=\scriptsize]{终稿文本回填} (3.2,-2.6);

        \node[note] at (1.6,-3.1) {\textit{路径B：文本（或语音终稿汇合）}};
        \draw[msg] (0,-3.5) -- node[above,font=\scriptsize]{文本输入} (3.2,-3.5);

        \draw[msg] (3.2,-4.1) -- node[above,font=\scriptsize]{\texttt{/ws/chat}连接（session,package,user\_id）} (12.0,-4.1);
        \draw[msg] (12.0,-4.5) -- node[above,font=\scriptsize]{加载catalog+TTS配置} (21.0,-4.5);
        \draw[retmsg] (21.0,-4.9) -- node[above,font=\scriptsize]{catalog数据} (12.0,-4.9);
        \draw[retmsg] (12.0,-5.3) -- node[above,font=\scriptsize]{\texttt{type:catalog}} (3.2,-5.3);

        \draw[msg] (3.2,-5.9) -- node[above,font=\scriptsize]{\texttt{\{"message":"...","scene\_location":"..."\}}} (12.0,-5.9);

        \draw[msg] (12.0,-6.5) -- node[above,font=\scriptsize]{读分层记忆+人设+用户画像} (21.0,-6.5);
        \draw[retmsg] (21.0,-6.9) -- node[above,font=\scriptsize]{瞬时List/短期List/长期String} (12.0,-6.9);

        \draw[msg] (12.0,-7.5) -- node[above,font=\scriptsize]{动作决策（非流式，含人设+瞬时记忆）} (15.3,-7.5);
        \draw[retmsg] (15.3,-7.9) -- node[above,font=\scriptsize]{\texttt{expression/motion} JSON} (12.0,-7.9);

        \draw[msg] (12.0,-8.5) -- node[above,font=\scriptsize]{聊天补全 stream=True} (15.3,-8.5);
        \draw[retmsg] (15.3,-8.9) -- node[above,font=\scriptsize]{token增量流} (12.0,-8.9);

        \draw[retmsg] (12.0,-9.5) -- node[above,font=\scriptsize]{\texttt{chunk}（首帧含expr/motion + content增量）} (3.2,-9.5);
        \node[note,anchor=west] at (3.3,-9.9) {\scriptsize setExpression/startMotion驱动Live2D};

        \node[note,anchor=east] at (18.1,-10.6) {\scriptsize 按句末标点攒批（默认N=4）并行worker合成};
        \draw[msg] (12.0,-10.9) -- node[above,font=\scriptsize]{TTS请求（文本段）} (18.2,-10.9);
        \draw[retmsg] (18.2,-11.4) -- node[above,font=\scriptsize]{WAV音频（保序）} (12.0,-11.4);

        \draw[retmsg] (12.0,-12.0) -- node[above,font=\scriptsize]{\texttt{chunk\_audio}\{index,size,content\}} (3.2,-12.0);
        \draw[retmsg,ultra thick] (12.0,-12.5) -- node[above,font=\scriptsize]{二进制WAV帧（index保序）} (3.2,-12.5);
        \node[note,anchor=west] at (3.3,-12.9) {\scriptsize enqueueAudioAndPlay；RMS$\rightarrow$口型同步};

        \draw[retmsg] (12.0,-13.6) -- node[above,font=\scriptsize]{\texttt{type:done}} (3.2,-13.6);

        \draw[msg] (12.0,-14.3) -- node[above,font=\scriptsize]{LPUSH瞬时/短期List；写chat\_session} (21.0,-14.3);
        \node[note,anchor=west] at (12.1,-14.8) {\scriptsize 每N轮或断连时异步调用remind\_extract抽取关怀触发器};
        \node[note,anchor=west] at (12.1,-15.3) {\scriptsize 每N轮或断连时异步刷新user\_profile};

        \foreach \x in {0,3.2,6.2,9.0,12.0,15.3,18.2,21.0}
            \fill (\x,-16.5) circle (2.5pt);
    \end{tikzpicture}
    }
    \caption{主对话端到端时序图（键盘与语音双入口汇合；朗读经同一\textbf{对话WebSocket接口}合并流下发；轮末触发记忆写回与关怀抽取）}
    \label{fig:seq_main_dialogue}
\end{figure}

% ====================================================
% 第四章 核心模块设计与实现
% ====================================================
\chapter{核心模块设计与实现}

本章按\textbf{浏览器呈现层→WebSocket协议层→FastAPI编排层→持久化与REST层}自顶向下展开，所有文件路径与环境变量默认值均以\texttt{Demo/}（前端）、\texttt{PY/}（后端）代码为准。\textbf{本章所述能力均面向已登录会话，不提供访客入口。}

\section{前端页面、Live2D 与交互实现}

\subsection{布局与登录态}

前端以Vite驱动模块化源码，\texttt{index.html}为单页入口。主界面采用\textbf{左会话、右角色}双栏：左侧集中历史气泡（\texttt{\#chat-list}）、流式字幕区、文本输入与语音按钮；右侧为WebGL画布与工具条。已登录用户在三步内完成``输入→等待流式反馈→同步听取朗读''。入口脚本在\texttt{DOMContentLoaded}后校验\texttt{localStorage}中的\texttt{live2d\_info}，缺失则重定向至\texttt{Login.html}。

\subsection{Live2D 组件层级}

图~\ref{fig:live2d_components}以UML组件图风格展示前端Live2D运行时四级层次：\texttt{LAppDelegate}（主代理，单例）→\texttt{LAppSubdelegate}（子代理/画布管理）→\texttt{LAppLive2DManager}（模型调度）→\texttt{LAppModel}（单模型实例）。\texttt{ws.js}经\texttt{applyChatLive2dActions}委托入口驱动表情与动作。

\begin{figure}[htbp]
    \centering
    \resizebox{0.88\textwidth}{!}{
    \begin{tikzpicture}[
        font=\small,
        comp/.style={rectangle, draw=black!70, thick, rounded corners=3pt, minimum width=4.0cm, minimum height=0.7cm, align=center, inner sep=5pt},
        agg/.style={{Diamond[open,length=3mm]}-{Latex[length=2.5mm]}, thick},
        dep/.style={-{Latex[length=2.5mm]}, thick, dashed},
        use/.style={-{Latex[length=2.5mm]}, thick, blue!70},
    ]
        \node[comp, fill=blue!10, minimum width=7.5cm] (delegate) at (5.0, 0)
            {\texttt{LAppDelegate}（单例）\\[-2pt]
             \scriptsize 初始化Cubism SDK；管理子代理集合；\\[-2pt]
             \scriptsize 注册全局指针/resize事件；帧循环广播；\texttt{applyChatLive2dActions}委托入口};
        \node[comp, fill=green!8, minimum width=6.5cm] (subdelegate) at (5.0, -3.0)
            {\texttt{LAppSubdelegate}\\[-2pt]
             \scriptsize 持有\texttt{<canvas>}+WebGL上下文；纹理管理；\\[-2pt]
             \scriptsize 画布尺寸/ResizeObserver；指针命中与拖拽};
        \node[comp, fill=orange!8, minimum width=6.0cm] (manager) at (5.0, -6.0)
            {\texttt{LAppLive2DManager}\\[-2pt]
             \scriptsize 维护\texttt{\_models[]}向量；切换ModelDir场景；\\[-2pt]
             \scriptsize 转发drag/tap/\texttt{applyChatLive2dActions}};
        \node[comp, fill=yellow!8, minimum width=5.5cm] (model) at (5.0, -9.0)
            {\texttt{LAppModel}\\[-2pt]
             \scriptsize 加载\texttt{.model3.json}、表情、物理、动作；\\[-2pt]
             \scriptsize \texttt{setExpression(id)}；\texttt{startMotionByBasename}};
        \node[comp, fill=gray!8, minimum width=3.0cm] (define) at (0.0, -9.0)
            {\texttt{lappdefine.js}\\[-2pt]\scriptsize ModelDir/资源根路径/常量};
        \node[comp, fill=purple!8, minimum width=4.5cm] (wsjs) at (11.5, -4.5)
            {\texttt{ws.js / main.js}\\[-2pt]
             \scriptsize 接收chunk/chunk\_audio；\\[-2pt]
             \scriptsize 解析expression/motion；\\[-2pt]
             \scriptsize Web Audio RMS→口型驱动};

        \draw[agg] (subdelegate.north) -- node[right,font=\scriptsize]{1..*} (delegate.south);
        \draw[agg] (manager.north) -- node[right,font=\scriptsize]{1} (subdelegate.south);
        \draw[agg] (model.north) -- node[right,font=\scriptsize]{1..*} (manager.south);
        \draw[dep] (define.east) -- node[above,font=\scriptsize]{使用} (model.west);
        \draw[use] (wsjs.west |- delegate.east) -- node[above,font=\scriptsize]{\texttt{applyChatLive2dActions}} (delegate.east);
        \draw[use,dashed] (wsjs.south) -- ++(0,-0.5) -- node[right,font=\scriptsize]{setExpr/startMotion} (model.east);
    \end{tikzpicture}
    }
    \caption{前端Live2D组件层级UML图（\texttt{LAppDelegate}管理多个\texttt{LAppSubdelegate}，\texttt{ws.js}经委托入口驱动表情与动作）}
    \label{fig:live2d_components}
\end{figure}

\subsection{远程资源与预签名加载}

私有MinIO桶场景下，\textbf{已登录用户}在\texttt{live2d\_info}校验通过后，由\texttt{main.js}执行：通过\texttt{GET /api/live2d-model-assets}拉取当前包资产行，对含\texttt{object\_key}的行调用\texttt{download-url}换取限时URL（默认\texttt{expires\_in=86400\,s}），最后将\texttt{relative\_path→可fetch URL}写入远程manifest供\texttt{LAppModel}拉取。并发解析默认并行上限24（\texttt{\_PRESIGN\_CONCURRENCY}），以平衡TLS握手与总等待。

\subsection{语音输入与对话入口合一}

语音链路单独维护至\texttt{/ws/asr}。浏览器以\texttt{MediaDevices.getUserMedia()}采集麦克风，PCM一般为16\,kHz/16\,bit/单声道，经二进制WebSocket帧上行；服务端\texttt{asr\_ws.py}与DashScope Fun-ASR建立实时会话并将事件回写。浏览器在收到终稿后填入输入框并调用与键盘相同的\texttt{sendChatMessage}，从而\textbf{共用同一对话状态机}。

\section{WebSocket 连接与流式协议}

\subsection{地址推导与查询参数语义}

\texttt{wsConfig.js}优先读取\texttt{VITE\_WS\_BASE}；否则由\texttt{VITE\_API\_BASE}推导\texttt{ws}/\texttt{wss}与主机端口；二者皆空时回退\texttt{ws://localhost:8000}。对话URL由\texttt{getChatWsUrl()}拼接三类查询参数：\texttt{session}（标签页级会话键，存\texttt{sessionStorage}，与MySQL \texttt{chat\_session.session\_key}对齐）、\texttt{package}（与\texttt{public/Resources/<package>}目录名一致）、\texttt{user\_id}（与已登录账户一致，供后端加载绑定的音色参考、人设等）。

\subsection{下行帧类型与朗读捆绑格式}

表~\ref{tab:ws_chat_frames}归纳\texttt{/ws/chat}上主要JSON类型字段时间顺序。二进制帧规则：仅当紧前一条JSON为\texttt{chunk\_audio}时，下一条二进制才被视作WAV；否则丢弃并告警。\texttt{ws.binaryType="arraybuffer"}，解码后进入Web Audio调度链，按\texttt{index}保序。

\begin{table}[htbp]
    \centering
    \caption{\texttt{/ws/chat}主要JSON类型与前端处理要点}
    \label{tab:ws_chat_frames}
    \small
    \setlength{\tabcolsep}{3.5pt}
    \begin{tabular}{|p{2.4cm}|p{5.4cm}|p{5.6cm}|}
        \hline
        \textbf{type} & \textbf{载荷要点} & \textbf{浏览器侧行为} \\
        \hline
        \texttt{catalog} & 包键、表情/动作名及路径列表 & 控制台诊断；约束LLM输出范围 \\
        \hline
        \texttt{chunk} & 文本增量\texttt{content}；首包可带\texttt{expression}/\texttt{motion} & 拼接字幕与气泡；驱动Live2D \\
        \hline
        \texttt{chunk\_audio} & \texttt{index}、\texttt{size}、对应文本段；可带头段动作字段 & 记录\texttt{pendingAudioMeta}；紧接二进制WAV \\
        \hline
        \texttt{remind\_trigger} & \texttt{trigger\_content}（情景）、\texttt{delivery\_message}（生成台词）等 & 气泡正文取\texttt{delivery\_message}；后可跟\texttt{chunk\_audio}+WAV朗读 \\
        \hline
        \texttt{done} & 本轮结束 & 合并完整气泡；收口型 \\
        \hline
        \texttt{error} & \texttt{message} & 追加错误提示；解除等待状态 \\
        \hline
    \end{tabular}
\end{table}

\subsection{断线重连、模型切换与消息打断}

客户端在\texttt{onclose}中以常量\texttt{RECONNECT\_MS=3000}定义的固定间隔（约3\,s）调度重连。切换模型时主动\texttt{close}并递增\texttt{skipAutoReconnectCount}，避免\texttt{onclose}触发的自动重连与手动换包重连叠加。\texttt{onmessage}首行比对\texttt{event.target===textSocket}，丢弃旧socket上迟到的JSON/二进制，防止换包或打断后字幕、音频错位入队。重连后\texttt{flush\_pending\_reminders\_for\_connection}会将该用户离线期间到期的关怀触发器补发。

\section{服务端对话编排}

\subsection{应用生命周期}

\texttt{PY/main.py}使用FastAPI \texttt{lifespan}：启动阶段调用\texttt{init\_catalog()}预热表情动作索引；\texttt{await start\_long\_memory\_consolidator()}挂载asyncio后台循环；\texttt{await start\_remind\_trigger\_scheduler()}启动关怀触发器常驻扫描任务。关闭阶段对称停止上述任务。CORS默认放行常见Vite开发端口，并允许通过环境变量\texttt{CORS\_ORIGINS}追加生产前端源。

\subsection{一轮对话内的编排流水线}

每收到非空\texttt{message}，编排顺序概括为：\textbf{(1)}记忆与人设拼装，按\texttt{memory\_layers.py}约定读取Redis列表与字符串、必要时回填MySQL长期概要，拼接system/user消息；\textbf{(2)}动作决策，在线程池中对\texttt{OLLAMA\_ACTION\_MODEL}做一次非流式补全，解析\texttt{expression}/\texttt{motion}，经集合校验与回退后附着于首段文本chunk；\textbf{(3)}聊天生成，对聊天模型启用\texttt{stream=True}，通过\texttt{asyncio.Queue}与事件循环线程桥接；\textbf{(4)}朗读，按句末标点攒批（默认N=4），分段并行请求TTS，合并流模式下以\texttt{chunk\_audio}+WAV保序下发；\textbf{(5)}收尾，发送\texttt{done}，持久化本轮对话至MySQL，更新Redis瞬时/短期列表，异步触发\texttt{remind\_extract}关怀抽取与用户画像刷新。

\subsection{TTS 切段保序流水线}

图~\ref{fig:tts_pipeline}展示服务端将聊天流切段、并行合成、保序下发的完整流水线，对应\texttt{wschat.py}中\texttt{sentence\_queue}、\texttt{tts\_worker}、\texttt{\_tts\_flush\_ordered}三段实现。

\begin{figure}[htbp]
    \centering
    \resizebox{0.9\textwidth}{!}{
    \begin{tikzpicture}[
        font=\small,
        box/.style={rectangle, draw=black!70, thick, rounded corners=3pt, minimum width=3.2cm, minimum height=0.75cm, align=center, inner sep=5pt},
        diam/.style={diamond, draw=black!70, thick, aspect=2.5, minimum width=3.8cm, minimum height=0.8cm, align=center, inner sep=2pt},
        arr/.style={-{Latex[length=2.5mm]}, thick},
        darr/.style={-{Latex[length=2mm]}, thick, dashed},
    ]
        \node[box, fill=blue!8] (tok) at (0,0) {Ollama token\\[-2pt]写入缓冲区};
        \node[diam, fill=orange!10, below=0.5cm of tok] (chk) {句末标点计数\\[-2pt]$\geq N$或token上限？};
        \node[box, fill=gray!8, right=1.2cm of chk] (cont) {继续缓冲};
        \node[box, fill=green!10, below=0.8cm of chk] (enq) {index自增\\[-2pt]文本段入sentence\_queue};
        \node[box, fill=blue!8, below=0.7cm of enq] (worker) {tts\_worker取\\[-2pt](index, text)};
        \node[box, fill=blue!8, below=0.5cm of worker] (synth) {asyncio.to\_thread\\[-2pt]调用TTS合成};
        \node[box, fill=green!10, below=0.5cm of synth] (store) {tts\_completed[index]\\[-2pt]$\leftarrow$ (text, WAV)};
        \node[box, fill=orange!8, below=0.6cm of store] (flush) {调用\_tts\_flush\_ordered()};
        \node[diam, fill=orange!10, below=0.5cm of flush] (rdy) {tts\_next\_send\\[-2pt]已就绪？};
        \node[box, fill=gray!8, right=1.2cm of rdy] (wait) {等待更小\\[-2pt]序号就绪};
        \node[box, fill=green!10, below=0.8cm of rdy] (send) {send\_json(chunk\_audio)\\[-2pt]+ send\_bytes(WAV)};
        \node[box, fill=blue!8, below=0.5cm of send] (inc) {tts\_next\_send += 1};

        \draw[arr] (tok) -- (chk);
        \draw[arr] (chk) -- node[left,font=\scriptsize]{是} (enq);
        \draw[darr] (chk) -- node[above,font=\scriptsize]{否} (cont);
        \draw[darr] (cont.south) |- ++(0,-0.3) -| (tok.east);
        \draw[arr] (enq) -- (worker);
        \draw[arr] (worker) -- (synth);
        \draw[arr] (synth) -- (store);
        \draw[arr] (store) -- (flush);
        \draw[arr] (flush) -- (rdy);
        \draw[arr] (rdy) -- node[left,font=\scriptsize]{是} (send);
        \draw[darr] (rdy) -- node[above,font=\scriptsize]{否} (wait);
        \draw[arr] (send) -- (inc);
        \draw[darr] (inc.east) -- ++(0.5,0) |- (flush.east);
    \end{tikzpicture}
    }
    \caption{TTS切段、并行合成与保序下发流程（对应\texttt{sentence\_queue}、\texttt{tts\_worker}、\texttt{\_tts\_flush\_ordered}）}
    \label{fig:tts_pipeline}
\end{figure}

\section{数据持久化、REST API 与资源管理}

\subsection{MySQL 八表结构}

模式定义于\texttt{schema.sql}（utf8mb4/InnoDB），图~\ref{fig:db_er}以简化ER图展示八张核心表的主要字段与关联关系。

\begin{figure}[htbp]
    \centering
    \includegraphics[width=\textwidth]{fig_er_diagram.png}
    \caption{MySQL八张核心表简化ER关系图（黄色底色为主键；蓝色箭头为外键约束）}
    \label{fig:db_er}
\end{figure}

\subsection{REST 路由分组与统一响应封装}

\texttt{http\_api.py}共注册\textbf{58}条\texttt{/api}路由。自定义\texttt{UnifiedResponseRoute}在路由处理层封装响应：成功时统一\texttt{code=0}、\texttt{message=ok}、\texttt{data}为业务载荷；失败时\texttt{code}取HTTP状态码、\texttt{message}为错误描述、\texttt{data=null}，便于前端单一解析逻辑。功能分组涵盖：用户与画像、\texttt{chat-sessions}、\texttt{long-memories}（含\texttt{consolidate-now}）、\texttt{personas}（含\texttt{expand-from-hints}）、\texttt{live2d-model-assets}（含\texttt{upload-zip}、\texttt{download-url}）、\texttt{live2d-tts-refers}、\texttt{remind-triggers}（含\texttt{scan-now}、\texttt{pending-scan}）、\texttt{background-images}。

\subsection{ZIP 入库与删除一致性}

\texttt{POST /api/live2d-model-assets/upload-zip}解压校验、推断包名、批量写MinIO与\texttt{live2d\_model\_asset}；删除包时移除对象与索引行，并清理依附\texttt{persona}及Redis缓存，保证再次连接\texttt{/ws/chat}时catalog与存储一致。

% ====================================================
% 第五章 分层记忆系统设计与实现
% ====================================================
\chapter{分层记忆系统设计与实现}

\section{理论参照}

\subsection{阿特金森--谢弗林多层存储模型}

认知心理学中的\textbf{阿特金森--谢弗林多层存储模型}（multi-store model）由Atkinson与Shiffrin系统阐述\cite{Atkinson1968}，核心图景是：记忆并非单一扁平仓库，而是在时间上串联的多级结构。外界信息先进入极短暂的感觉登记（sensory register），经注意进入容量有限的短时存储（short-term store，STS），继而在复述、编码等控制过程作用下转入保持更久的长时存储（long-term store，LTS），以供日后检索与利用。

\textbf{本文立场（类比而非仿真）：}上文关于三级结构的概括用于读者定位本文瞬时/短期/长期三联设计的叙述来源，绝不等于在计算机里复制感官生理过程或复述短时容量等心理学定量结论。表~\ref{tab:memory_analogy}给出工程对应与类比边界说明。

\begin{table}[htbp]
    \centering
    \caption{阿特金森--谢弗林三级结构与本文三层记忆的叙述类比（非等价映射）}
    \label{tab:memory_analogy}
    \footnotesize
    \setlength{\tabcolsep}{3pt}
    \begin{tabular}{p{3.35cm}p{5.05cm}p{5.25cm}}
        \hline
        \textbf{心理学层级（简述）} & \textbf{本文工程对应（机制要点）} & \textbf{类比边界说明} \\
        \hline
        感觉登记：极短保持的感觉缓冲 & （不对齐独立模块）人机采集层PCM、ASR流式中间结果等 & 仅说明``门户缓冲''直觉；无感官生理学仿真 \\
        \hline
        短时存储：容量/时间受限，复述有助于转入长时 & 瞬时List（最近$K$轮）+短期List（挤出后的规则条目） & 容量由常量与TTL约束；``复述''对应规则压缩，而非心理复述 \\
        \hline
        长时存储：相对持久、可提取 & \texttt{long\_memory}周期概要（MySQL）+Redis长期字符串缓存 & 非完整语义记忆理论实现；强调可注入、可审计、可失效（TTL） \\
        \hline
    \end{tabular}
\end{table}

\subsection{MemGPT 的工程类比}

MemGPT的核心直觉是：大模型单次能读进去的字数有限，不可能把全部历史原封不动塞进每次请求；因此把\textbf{当下真正拼进prompt的一小段}与\textbf{篇幅更长、存放在外部的记录}分开——后者落在数据库或向量库等外部存储里，对话进行中再按需读取、拼接或摘要写回，从而在不要求单次prompt无限变长的前提下延续多轮一致性\cite{MemGPT2023}。

与MemGPT相比，本文要解决的是同一件实事：上下文长度有限，却又要利用更早的对话与人设。\textbf{写法上同样拆成两块：}第一块是本轮拼进请求里的文字——主要是最近若干轮原文及短期压缩条目；第二块是平时放在Redis列表、字符串或MySQL字段里的较长记录，只在后端代码的固定步骤里读取或更新。

\textbf{本文刻意不做或少做的MemGPT要素：}一是由模型自主决定何时读写库（本文把何时写Redis、何时读MySQL概要写死在编排代码的常见路径里）；二是论文中与操作系统类比的复杂调度（本文用的是可在配置里改、可在日志里核对的Redis List操作与TTL）；三是检索形态（本文长期层以结构化表字段+文本概要为主，优先保证工程落地与数据所有权清晰）。

\section{三层架构设计}\label{sec:memgptrelation}

图~\ref{fig:memory_arch}展示三层记忆的存储介质、数据形态与注入方式。

\begin{figure}[htbp]
    \centering
    \resizebox{0.95\textwidth}{!}{
    \begin{tikzpicture}[
        font=\small,
        layer/.style={rectangle, draw=black!70, thick, rounded corners=4pt, minimum width=13cm, inner sep=10pt},
        sublabel/.style={font=\bfseries\small, anchor=north west},
        detail/.style={font=\scriptsize, anchor=north west},
        arr/.style={-{Latex[length=3mm]}, thick},
    ]
        \node[layer, fill=blue!8] (instant) at (0,0) {
            \begin{minipage}{12.6cm}
            \textbf{瞬时记忆层（Instant Memory）}\\[2pt]
            \scriptsize 存储介质：Redis List，键前缀默认\texttt{moment}，空闲TTL默认3600\,s\\
            \scriptsize 数据形态：JSON元素\texttt{\{"u":"用户发言","a":"助手回复","ts":"时间戳"\}}，每元素一轮\\
            \scriptsize 容量策略：\texttt{LPUSH}新轮入头，\texttt{LTRIM}保留最近$K$轮（默认$K=5$）\\
            \scriptsize 注入方式：直接展开为\texttt{messages}中时间顺序的\texttt{user/assistant}多轮对话\\
            \scriptsize 动作LLM：独立读取同一瞬时List，经\texttt{\_format\_instant\_turns\_for\_action\_llm}截断排版后注入\texttt{user}侧
            \end{minipage}
        };
        \node[layer, fill=green!8, below=0.4cm of instant] (short) {
            \begin{minipage}{12.6cm}
            \textbf{短期记忆层（Short-term Memory）}\\[2pt]
            \scriptsize 存储介质：Redis List，键前缀默认\texttt{short}，单键TTL默认86400\,s\\
            \scriptsize 数据形态：两类条目——规则压缩条（\texttt{type:rule}）与摘要条（\texttt{type:summary}，结构预留）\\
            \scriptsize 容量策略：最多20条（\texttt{SHORT\_TERM\_MAX\_ENTRIES}）；拼入system的字符上限默认4000\\
            \scriptsize 写入时机：瞬时List溢出时，被\texttt{LTRIM}挤出的最旧轮经\texttt{rule\_compact\_turn}规则压缩后\texttt{LPUSH}入本层\\
            \scriptsize 注入方式：\texttt{LRANGE}后从新到旧格式化为纯文本，附加在包级system正文之后的``短期记忆''段落
            \end{minipage}
        };
        \node[layer, fill=orange!8, below=0.4cm of short] (longm) {
            \begin{minipage}{12.6cm}
            \textbf{长期记忆层（Long-term Memory）}\\[2pt]
            \scriptsize 存储介质：MySQL \texttt{long\_memory.period\_overview}（TEXT）+Redis String缓存（键前缀\texttt{long}，TTL默认604800\,s）\\
            \scriptsize 数据形态：\texttt{period\_overview}存周期性Ollama生成的口述式概要；新概要以分隔线追加\\
            \scriptsize 更新时机：后台\texttt{long\_memory\_consolidator}每约24\,h扫描MySQL \texttt{chat\_session}时间窗并生成；亦可\texttt{POST /api/long-memories/consolidate-now}即时触发\\
            \scriptsize 注入方式：读Redis String（未命中则回填MySQL），注入上限默认2000字符，附加在system的``长期记忆''段落
            \end{minipage}
        };

        \draw[arr] (instant.south) -- node[right,font=\scriptsize]{溢出轮次经规则压缩} (short.north);
        \draw[arr] (short.south) -- node[right,font=\scriptsize]{周期概要（后台任务）} (longm.north);
    \end{tikzpicture}
    }
    \caption{三层记忆架构：瞬时层直接展开为对话messages；短期层经规则压缩注入system；长期层由后台任务周期生成并缓存于Redis}
    \label{fig:memory_arch}
\end{figure}

\section{瞬时记忆：Redis List 读写机制}

\subsection{JSON 元素格式与 LPUSH/LTRIM 策略}

每轮对话结束后，\texttt{\_append\_turn\_memory\_layers}执行下列写入：

\begin{verbatim}
算法 瞬时记忆写入（每轮对话结束后执行）
输入: user_id, package_key, 本轮用户文本 U, 助手文本 A
1. 构造 JSON 轮次 R ← {"u":U, "a":A, "ts":now()}
2. LPUSH instant_key(user, pkg), R
3. LTRIM instant_key 保留最近 K 轮         // 默认 K=5
4. 若本轮挤出最旧轮 E:
5.     E' ← rule_compact_turn(E)          // 截断+敷衍过滤
6.     LPUSH short_key, E'
7.     对 short_key 执行条数裁剪与 TTL 续期
输出: 更新后的 Redis 状态
\end{verbatim}

\subsection{注入主对话 messages 的方式}

\texttt{\_build\_memory\_for\_model}通过\texttt{read\_instant\_turns\_chronological}读取瞬时List后，按时间正序展开为多条\texttt{user}/\texttt{assistant}消息，插入当前轮用户消息之前。这一方式使模型能直接``看见''最近$K$轮完整对话原文，无需经过任何压缩损耗。

\subsection{动作 LLM 对瞬时记忆的独立使用}

动作/表情选型模型为\textbf{单独一次}非流式调用，其\texttt{user}侧由\texttt{\_action\_llm\_user\_content\_sync}拼装：在$\texttt{user\_id}\geq 1$且能读到\texttt{Persona}时追加人设块；Redis可用时追加瞬时记忆格式化正文；末段为本轮用户输入。\textbf{不注入}短期或长期记忆，\textbf{亦不注入}用户画像块。选型语义明确：\texttt{expression}/\texttt{motion}表示\textbf{Live2D虚拟角色（助手）}应切换到的资源，不是真人用户的表情或动作。

\section{短期记忆：规则压缩与 system 拼装}

\subsection{规则压缩算法}

被\texttt{LTRIM}挤出的最旧轮次经\texttt{rule\_compact\_turn}处理：对用户句与助手句分别做长度截断；若助手回复为敷衍应答（极短且无实质内容），则清空\texttt{ai\_response}字段。压缩后条目形态为：

\begin{verbatim}
{
  "type": "rule",
  "time": "2026-05-08T10:30:00+00:00",
  "user_question": "（截断后的用户句）",
  "ai_response": "（截断后的助手句，敷衍时为空）"
}
\end{verbatim}

\subsection{摘要条目结构预留}

短期List亦可存放\texttt{type:summary}条目：默认由\texttt{utils/short\_term\_summary.maybe\_push\_short\_term\_summary\_after\_turn}在\texttt{/ws/chat}每轮写Redis后按独立计数键每\texttt{N}轮（环境变量控制，默认与画像刷新同为5）调用Ollama生成摘要并调用\texttt{push\_summary\_entry}；亦可由业务侧直接调用\texttt{push\_summary\_entry}。结构形态为\texttt{\{"type":"summary","time":...,"text":"摘要正文"\}}。写入时可从短期List中移除\texttt{type==rule}且时间落在\texttt{prune\_rule\_turn\_times}内的条目，避免摘要与规则条重复占用上下文。\texttt{SHORT\_TERM\_SUMMARY\_ENABLED=1}开启内置周期摘要（默认为关闭以免额外占用推理）。

\subsection{system 拼装与字符上限控制}

\texttt{\_build\_memory\_for\_model}生成的\texttt{system}首条基底由\texttt{\_chat\_system\_prompt\_for\_session}产生：全局指令→包级人设（\texttt{character\_desc}/\texttt{tone\_style}）→可选叙事语境块（\texttt{scene\_location}/\texttt{scene\_time}）→可选用户画像块。随后按需拼接长期记忆段落（字符上限默认2000）与短期记忆段落（字符上限默认4000），按从新到旧顺序格式化后附加。

\section{长期记忆：周期概要后台任务}

\subsection{long\_memory\_consolidator 调度机制}

后台任务\texttt{long\_memory\_consolidator.py}由\texttt{PY/main.py}的\texttt{lifespan}调用\texttt{start\_long\_memory\_consolidator}挂载，默认每约\textbf{24小时}扫描一轮候选用户-包对（通过\texttt{list\_candidates\_for\_consolidation}施加最短24小时间隔限制），执行\texttt{consolidate\_one}：

\begin{enumerate}
    \item 读取MySQL \texttt{chat\_session}最近24小时时间窗内会话行（\texttt{list\_for\_long\_memory\_window}）；
    \item 将会话规则合并为口述式叙述，调用Ollama生成\texttt{period\_overview}（可做严格重试与\texttt{finalize}修补）；
    \item \texttt{upsert\_by\_user\_pkg}写入MySQL \texttt{long\_memory}，新概要以分隔线追加至已有正文；
    \item 调用\texttt{write\_long\_memory\_text}刷新Redis长期String缓存。
\end{enumerate}

\subsection{即时周期概要更新接口}

\texttt{POST /api/long-memories/consolidate-now}对指定\texttt{user\_id+package\_key}调用同一套\texttt{consolidate\_one}，在\texttt{asyncio.to\_thread}中执行以免阻塞事件循环；\textbf{不受}后台候选筛选里24小时最短间隔限制，便于开发期手动触发或前端工具栏``总结长期记忆''功能调用。

\section{用户画像：user\_profile 的生成与注入}

\subsection{触发条件与 LLM 合并机制}

用户画像（\texttt{user\_profile}）每\texttt{user\_id}一行，字段包含\texttt{user\_tags}、\texttt{emotion\_state}、\texttt{preferences}、\texttt{trouble\_events}。刷新由\texttt{user\_profile\_refresh.py}负责，触发条件有两类：一是\texttt{maybe\_refresh\_user\_profile\_after\_turn}，Redis计数键能被$N$整除时（默认$N=5$）触发；二是\texttt{refresh\_user\_profile\_on\_disconnect}，\texttt{/ws/chat}的\texttt{finally}中无论轮次计数均触发一次。刷新时以旧画像字段格式化为正文并拼接当前包下瞬时记忆多轮对话，调用LLM生成结构化字段后\texttt{upsert}至MySQL。

\subsection{画像与人设的分列设计原则}

系统中用户画像描述的是\textbf{真人用户}的标签与偏好，用于共情与贴合用户应答；\textbf{人设}（persona）描述的是\textbf{虚拟角色自身}的人格与表达方式。二者之所以\textbf{严格分列}，来自以下几点设计考量：

\textbf{其一，个性化体验的完整性。}真人用户的偏好（如音乐口味、情绪状态）与角色的语气风格（如温柔知性、幽默活泼）是两个独立维度；若二者混入同一描述段落，LLM往往难以分清``谁是谁''，导致回复风格混乱或角色``人格塌陷''。分列设计使个性化在用户侧和角色侧同时发力，而不相互干扰。此外，系统支持从后端随机背景图库中为每次会话加载不同场景背景，背景名称同样注入对话语境块，使角色在不同叙事场景中对应不同氛围，进一步丰富个性化层次。

\textbf{其二，防止角色退化为通用助手。}角色人设（含性格描述与语气风格）明确告知模型``这是一位有自己世界观的虚拟个体''；若将用户标签与角色特征堆砌在一起，模型容易退化为``知道很多用户信息的通用客服''，丧失角色独特性。分列后，人设段落专注于塑造角色人格，画像段落专注于提供用户背景参照，两条线各司其职。

\textbf{其三，为角色故事深化预留空间。}人设字段目前包含角色描述与语气风格，未来可在此基础上扩展角色成长轨迹、情感状态历史与世界观设定，逐步形成``数字角色故事线''，使每位用户与同一角色的互动随着时间积累而具有连贯性与专属性。分列结构使这一扩展路径清晰，无需重构记忆注入逻辑。

在实现上，用户画像注入主对话系统提示（可通过环境变量开关控制，默认开启）；人设同样注入系统提示，但二者分属不同段落，在关怀话术生成时亦被分别引用，保证``画像描述用户、人设描述角色''的语义边界在全链路中一致维持。

\section{分层记忆整体读写流程}

图~\ref{fig:memory_flow}展示服务端在一轮对话中读取并组装记忆上下文（轮首）以及轮末将本轮对话写回各层存储的完整流程。

\begin{figure}[htbp]
    \centering
    \resizebox{0.95\textwidth}{!}{
    \begin{tikzpicture}[
        font=\small,
        box/.style={rectangle, draw=black!70, thick, rounded corners=3pt, minimum width=3.5cm, minimum height=0.7cm, align=center, inner sep=4pt},
        redbox/.style={rectangle, draw=red!60, thick, rounded corners=3pt, fill=red!5, minimum width=3.5cm, minimum height=0.7cm, align=center, inner sep=4pt},
        sqlbox/.style={rectangle, draw=blue!60, thick, rounded corners=3pt, fill=blue!5, minimum width=3.5cm, minimum height=0.7cm, align=center, inner sep=4pt},
        arr/.style={-{Latex[length=2.5mm]}, thick},
        darr/.style={-{Latex[length=2mm]}, thick, dashed, gray!70},
    ]
        \node[box, fill=blue!8] (start) at (0,0) {收到用户消息};
        \node[redbox] (inst) at (-4.5,-2) {读瞬时List\\[-2pt]\texttt{moment:*}（最近K轮）};
        \node[redbox] (short) at (0,-2) {读短期List\\[-2pt]\texttt{short:*}（最多20条）};
        \node[redbox] (longr) at (4.5,-2) {读长期String\\[-2pt]\texttt{long:*}（周期概要）};
        \node[sqlbox] (longm) at (4.5,-3.5) {Redis未命中\\[-2pt]回填MySQL \texttt{long\_memory}};
        \node[box, fill=green!8] (assem) at (0,-5.0) {拼装system/user\\[-2pt]记忆上下文};
        \node[box, fill=orange!8, below=0.6cm of assem] (infer) {调用Ollama\\[-2pt]动作决策+对话生成};
        \node[box, fill=blue!8, below=0.6cm of infer] (done) {对话轮结束\\[-2pt]\texttt{type:done}};
        \node[redbox] (winst) at (-4.5,-10.5) {LPUSH瞬时List\\[-2pt]LTRIM保留K轮};
        \node[redbox] (wshort) at (0,-10.5) {被挤出轮次\\[-2pt]规则压缩入短期List};
        \node[sqlbox] (wmysql) at (4.5,-10.5) {持久化\\[-2pt]\texttt{chat\_session}};
        \node[sqlbox] (wlong) at (2.0,-12.2) {后台任务（周期）\\[-2pt]写\texttt{long\_memory.period\_overview}};
        \node[redbox] (wlongr) at (2.0,-13.5) {刷新Redis\\[-2pt]长期String缓存};

        \draw[arr] (start) -- ++(-4.5,-1.3) -- (inst);
        \draw[arr] (start) -- (short);
        \draw[arr] (start) -- ++(4.5,-1.3) -- (longr);
        \draw[darr] (longr) -- (longm);
        \draw[darr] (longm.south) |- ++(0,-0.3) -| (assem.east);
        \draw[arr] (inst) |- (assem.west);
        \draw[arr] (short) -- (assem);
        \draw[arr] (longr) |- (assem.east);
        \draw[arr] (assem) -- (infer);
        \draw[arr] (infer) -- (done);
        \draw[arr] (done) -- ++(-4.5,-1.3) -- (winst);
        \draw[arr] (done) -- (wshort);
        \draw[arr] (done) -- ++(4.5,-1.3) -- (wmysql);
        \draw[darr] (winst.south) |- ++(0,-0.5) -- (wshort.south);
        \draw[darr] (wmysql.south) |- (wlong);
        \draw[arr] (wlong) -- (wlongr);
    \end{tikzpicture}
    }
    \caption{分层记忆读写流程图（轮首由\texttt{\_build\_memory\_for\_model}从Redis/MySQL拼装；轮末由\texttt{\_append\_turn\_memory\_layers}写回；长期层由后台consolidator周期更新）}
    \label{fig:memory_flow}
\end{figure}

% ====================================================
% 第六章 主动关怀系统设计与实现
% ====================================================
\chapter{主动关怀系统设计与实现}

\section{设计动机：从被动对话到时间驱动关怀}

传统对话系统的交互模式是严格被动的：系统仅在用户发起消息后才产生响应，无法感知用户在对话间隙的状态变化，更无法在约定时刻主动触达用户。然而，真实的情感陪伴往往需要系统具备时间感知能力——在用户面临考试、生日或低谷期时主动给予关怀，而不是等待用户先开口。本章所述主动关怀系统正是为此而设计：通过对话内信息抽取、触发器持久化与时间驱动调度，实现从``被动应答''到``主动关怀''的跨越。

图~\ref{fig:care_loop}给出四步闭环的整体设计。

\begin{figure}[htbp]
    \centering
    \resizebox{0.82\textwidth}{!}{
    \begin{tikzpicture}[
        font=\small,
        stepbox/.style={rectangle, draw=black!70, thick, rounded corners=6pt, minimum width=3.5cm, minimum height=1.4cm, align=center, inner sep=6pt},
        arr/.style={-{Latex[length=3mm]}, ultra thick, gray!60},
    ]
        \node[stepbox, fill=blue!10]  (s1) at (0,  0)    {\textbf{①记忆存储}\\[-2pt]\scriptsize 对话内容沉淀于\\[-2pt]\scriptsize 瞬时/短期/长期\\[-2pt]\scriptsize 三级记忆与user\_profile};
        \node[stepbox, fill=orange!10](s2) at (6.5,0)    {\textbf{②信息抽取}\\[-2pt]\scriptsize remind\_extract\\[-2pt]\scriptsize 异步从对话中\\[-2pt]\scriptsize 抽取结构化提醒};
        \node[stepbox, fill=green!10] (s3) at (6.5,-4.5) {\textbf{③定时扫描}\\[-2pt]\scriptsize 后台调度器\\[-2pt]\scriptsize 周期扫描到期\\[-2pt]\scriptsize 触发器并认领};
        \node[stepbox, fill=purple!10](s4) at (0,-4.5)   {\textbf{④生成与推送}\\[-2pt]\scriptsize remind\_delivery\\[-2pt]\scriptsize LLM生成话术\\[-2pt]\scriptsize WebSocket投递};

        \draw[arr] (s1.east) -- node[above,font=\scriptsize]{对话落库后异步} (s2.west);
        \draw[arr] (s2.south) -- node[right,font=\scriptsize]{触发器入库} (s3.north);
        \draw[arr] (s3.west) -- node[above,font=\scriptsize]{到期认领} (s4.east);
        \draw[arr] (s4.north) -- node[left,font=\scriptsize]{用户接收关怀\\记忆更新} (s1.south);
    \end{tikzpicture}
    }
    \caption{主动关怀系统四步闭环设计（①记忆存储→②信息抽取→③定时扫描→④生成与推送，形成主动关怀循环）}
    \label{fig:care_loop}
\end{figure}

\section{对话内关怀信息抽取（remind\_extract）}

每轮对话\texttt{chat\_session}成功落库后，\texttt{wschat.py}异步调用\texttt{PY/utils/remind\_extract.py}，由LLM从本轮用户话+助手回复中抽取结构化提醒，写入\texttt{remind\_trigger}表，并绑定\textbf{刚插入行的\texttt{session\_id}}。

抽取结果写入\texttt{remind\_trigger}的字段语义如下：\texttt{trigger\_type}描述关怀类型（如``生日''``考试''``纪念日''``日常关怀''等）；\texttt{trigger\_time}为应触发的时刻；\texttt{trigger\_content}存\textbf{情景详细描述}（即表内对当时情景的总结）；\texttt{session\_id}关联产生该提醒的单轮\texttt{chat\_session}，供后续话术生成时读取关联对话上下文。\texttt{is\_triggered}初始为0，被调度器认领后置1。

实现上将\textbf{服务器当前时间}（进程本地\texttt{datetime.now()}快照）与\textbf{对话语境参考时间}（由本条\texttt{/ws/chat} JSON中的\texttt{scene\_time}经解析得到，解析失败则等同服务器时间）写入抽取模型的用户提示，要求其\textbf{自行推算}\texttt{trigger\_time}字符串（含相对短语到绝对时刻的换算）；提示词约束公历年默认与参考自然年一致，\textbf{除非}对话明确``明年''等。\textbf{不在服务端用正则覆盖模型给出时刻}，亦无``模型未产出条目时的兜底插入''。\texttt{\_parse\_trigger\_time}将模型字符串解析为\texttt{DATETIME}，允许粗粒度（如仅\texttt{YYYY-MM}、仅日期、或到小时），缺失分秒按约定补零或与\texttt{09:00:00}对齐；日历类无年份/早于当前的点仍经\texttt{\_roll\_trigger\_time\_to\_future}滚入未来，但对``已是未来却年份无理''等情况不做代码纠偏。

\section{触发器数据结构与 REST 接口}

\texttt{remind\_trigger}表是主动关怀的核心数据结构，完整字段见表~\ref{tab:remind_trigger_fields}。

\begin{table}[htbp]
    \centering
    \caption{\texttt{remind\_trigger}表字段说明}
    \label{tab:remind_trigger_fields}
    \small
    \setlength{\tabcolsep}{4pt}
    \begin{tabular}{|p{3.2cm}|p{2.5cm}|p{6.7cm}|}
        \hline
        \textbf{字段名} & \textbf{类型/约束} & \textbf{含义} \\
        \hline
        \texttt{trigger\_id} & BIGINT PK & 触发器唯一ID \\
        \hline
        \texttt{user\_id} & INT FK→user & 关联用户 \\
        \hline
        \texttt{session\_id} & BIGINT FK→chat\_session & 产生该提醒的对话轮次，用于话术生成时读取关联上下文 \\
        \hline
        \texttt{trigger\_type} & VARCHAR(30) & 触发类型（生日/考试/纪念日/日常关怀等） \\
        \hline
        \texttt{trigger\_time} & DATETIME & 应触发的时刻（到达此时刻且is\_triggered=0才纳入扫描） \\
        \hline
        \texttt{trigger\_content} & TEXT & 库内存情景详细描述（JSON 中与 REST/WebSocket 字段同名同义） \\
        \hline
        \texttt{is\_triggered} & TINYINT DEFAULT 0 & 0=未触发；1=已触发（认领后置1） \\
        \hline
        \texttt{create\_time} & DATETIME & 写入时间 \\
        \hline
    \end{tabular}
\end{table}

REST接口提供完整的CRUD与诊断能力：\texttt{GET /api/remind-triggers/pending-scan}拉取截至当前时刻仍未触发且已到期的记录列表；\texttt{GET /api/remind-triggers/pending-count}返回其数量；\texttt{POST /api/remind-triggers/scan-now}立即执行一轮扫描，响应\texttt{data}中含\texttt{pending\_fetched}、\texttt{delivered}、\texttt{released\_no\_ws}等统计字段，便于排查``无推送''原因是否为无到期行或用户未在线。

\section{后台调度器}

\subsection{常驻扫描机制}

\texttt{PY/main.py}的\texttt{lifespan}启动\texttt{PY/live2d\_db/remind\_trigger\_scheduler.py}，按间隔轮询到期行（默认\textbf{300秒}，可在\texttt{PY/.env}调小至最小5秒，便于联调）；扫描条件与\texttt{GET pending-scan}使用同一套仓储方法\texttt{list\_pending\_before}：即\texttt{is\_triggered=0}且\texttt{trigger\_time<=服务器当前时间}。

\subsection{「到期」与「扫描发现」两阶段语义（重要）}\label{sec:two_phase}

定时关怀\textbf{不是}在\texttt{trigger\_time}某一时刻由独立闹钟线程精确唤醒，而是同时满足以下两点才对用户可见：

\begin{enumerate}
    \item \textbf{到期（数据条件）}：表中该行\texttt{is\_triggered=0}，且\texttt{trigger\_time<=服务器当前时间}。若面试定在6月15日09:00，在此时间之前无论如何点击扫描，都会出现``当前无到期未触发记录''——这是预期行为。
    \item \textbf{发现（调度条件）}：在满足(1)之后，还须经过\textbf{一次扫描}：后台定时任务或\texttt{POST /api/remind-triggers/scan-now}。因此推送时刻是``到期之后的下一次扫描''，存在至多约一个扫描间隔的滞后（默认约5分钟；联调可缩短扫描间隔，或将测试数据的\texttt{trigger\_time}设为``几分钟前''再点立即扫描）。
\end{enumerate}

\textbf{离线用户}：无在线\texttt{/ws/chat}时不会写入已触发；记录保持未触发，待用户重新连接时由\texttt{flush\_pending\_reminders\_for\_connection}批量补发并认领。

\subsection{乐观认领机制}

调度器通过\texttt{claim\_pending\_trigger}对记录进行乐观认领：先将\texttt{is\_triggered}置1并加时间戳，再尝试投递；若投递失败（用户不在线）则释放认领（\texttt{is\_triggered}回置0），供下次扫描重试。这一机制保证同一触发器不会被并发重复推送。

\section{关怀话术生成（remind\_delivery）}\label{sec:delivery}

\subsection{输入拼装}

\texttt{PY/utils/remind\_delivery.py}在调度器成功认领后被调用，其输入拼装顺序（写入LLM的User消息）为：

\begin{enumerate}
    \item \textbf{人设块（可选）}：当前包的\texttt{character\_desc}/\texttt{tone\_style}，通过\texttt{\_persona\_block\_for\_package(user\_id, pkg)}读取；
    \item \textbf{用户画像块（可选）}：当\texttt{USER\_PROFILE\_IN\_CHAT}未关闭且\texttt{user\_profile}存在非空字段时，由\texttt{format\_profile\_for\_chat\_system}格式化为【用户画像】段落，便于生成关怀台词时参考用户背景；
    \item \textbf{当前对话瞬时记忆（可选）}：与同连接\texttt{/ws/chat}主对话相同\texttt{(user\_id, package\_key)}分桶，从Redis瞬时List按时间顺序读取最近多轮\texttt{user}/\texttt{assistant}原文；\textbf{不}拼接短期压缩层条目；
    \item \textbf{关怀类型}：\texttt{trigger\_type}；
    \item \textbf{情景详细描述}：库内\texttt{trigger\_content}（无则写``未提供''）；
    \item \textbf{关联对话原文}：通过\texttt{session\_id}读取\texttt{chat\_session}的\texttt{user\_input}/\texttt{ai\_reply}；无有效\texttt{session\_id}时写``未绑定''。
\end{enumerate}

LLM系统提示要求\textbf{重新撰写}对用户当场说出的一段话：以情景与瞬时对话为主轴，可呼应关联单轮与人设/画像；输出为1～3句口语化中文、不编造材料中不存在的事实、禁止JSON/Markdown格式；输出长度上限约800字符，超出截断加``…''。

\subsection{回退策略}

在\texttt{PY/.env}将``关怀话术是否走LLM''设为关闭时，不调用Ollama，直接走\texttt{\_fallback\_line}：有情景描述则截断约240字（过长加省略号），否则一句默认问候。此回退保证在LLM不可用时系统仍能正常投递关怀消息。

\section{WebSocket 推送与 TTS 朗读}

\subsection{type:remind\_trigger 帧格式}

认领成功后，\texttt{broadcast\_remind\_trigger\_to\_user}向该\texttt{user\_id}当前所有\texttt{/ws/chat}连接发送：

\begin{verbatim}
{
  "type": "remind_trigger",
  "trigger_id": <int>,
  "trigger_type": "<类型>",
  "trigger_content": "<库内情景详细描述>",
  "delivery_message": "<当场LLM生成的面向用户台词>",
  "trigger_time": "<ISO时间>"
}
\end{verbatim}

前端\texttt{ws.js}以\texttt{delivery\_message}作为气泡正文；以\texttt{appendChatMessage("remind", ...)}插入聊天列表，米色气泡\texttt{chat-item.remind}与普通对话区分。

\subsection{可选 MiMo 语音朗读}

若MiMo已在\texttt{PY/.env}完整配置且未关闭``关怀语音''开关，则在\texttt{type:remind\_trigger}帧之后，再按该连接\texttt{package}的参考音与人设导演下发\texttt{chunk\_audio}（含\texttt{remind\_audio:true}标记）+二进制WAV，朗读内容与JSON中\texttt{delivery\_message}一致，与合并流聊天同一代码路径。关闭开关则仅保留文字气泡，音频不下发。

\subsection{离线用户重连补发}

用户离线期间到期的触发器记录保持\texttt{is\_triggered=0}（认领失败后释放）。当用户重新建立\texttt{/ws/chat}连接时，\texttt{flush\_pending\_reminders\_for\_connection}被调用，以\texttt{list\_pending\_for\_user\_before(now)}读取所有属于该用户的到期未触发记录，依次走与常驻扫描相同的认领→话术生成→投递链路，确保用户不会因短暂离线而错过关怀消息。

\section{主动关怀端到端时序}

图~\ref{fig:care_seq}展示从对话落库到WebSocket推送的完整时序，涵盖抽取、调度、话术生成与推送四个阶段。

\begin{figure}[htbp]
    \centering
    \resizebox{\textwidth}{!}{
    \begin{tikzpicture}[
        font=\small,
        participant/.style={rectangle, draw=black!70, thick, minimum width=2.3cm, minimum height=0.6cm, align=center, fill=white},
        lifeline/.style={dashed, thick, gray!50},
        msg/.style={-{Latex[length=2mm]}, thick},
        retmsg/.style={-{Latex[length=2mm]}, thick, dashed},
        note/.style={rectangle, draw=black!35, fill=yellow!10, rounded corners=2pt, font=\scriptsize, inner sep=2.5pt, align=left},
    ]
        \node[participant] (dial) at (0,0)    {对话编排\\[-2pt]{\scriptsize /ws/chat}};
        \node[participant] (ext)  at (4.0,0)  {关怀抽取\\[-2pt]{\scriptsize remind\_extract}};
        \node[participant] (db)   at (8.0,0)  {MySQL\\[-2pt]{\scriptsize remind\_trigger}};
        \node[participant] (sched)at (12.0,0) {后台调度器\\[-2pt]{\scriptsize scheduler}};
        \node[participant] (del)  at (16.0,0) {话术生成\\[-2pt]{\scriptsize remind\_delivery}};
        \node[participant] (fe)   at (20.0,0) {浏览器\\[-2pt]{\scriptsize ws.js}};

        \foreach \x in {0,4.0,8.0,12.0,16.0,20.0}
            \draw[lifeline] (\x,-0.38) -- (\x,-14.0);

        % Phase 1: Extract
        \node[note] at (1.8,-0.9) {\textit{阶段①：对话抽取（每轮落库后异步）}};
        \draw[msg] (0,-1.4) -- node[above,font=\scriptsize]{对话落库成功} (4.0,-1.4);
        \draw[msg] (4.0,-2.0) -- node[above,font=\scriptsize]{LLM抽取结构化提醒} (8.0,-2.0);
        \draw[retmsg] (8.0,-2.6) -- node[above,font=\scriptsize]{写入remind\_trigger（绑定session\_id）} (4.0,-2.6);

        % Phase 2: Scheduling
        \node[note] at (9.8,-3.3) {\textit{阶段②：定时扫描（默认300s间隔）}};
        \draw[msg] (12.0,-3.8) -- node[above,font=\scriptsize]{list\_pending\_before(now)} (8.0,-3.8);
        \draw[retmsg] (8.0,-4.3) -- node[above,font=\scriptsize]{到期且is\_triggered=0的记录} (12.0,-4.3);
        \draw[msg] (12.0,-4.9) -- node[above,font=\scriptsize]{claim（is\_triggered←1）} (8.0,-4.9);

        % Phase 3: Delivery
        \node[note] at (13.8,-5.6) {\textit{阶段③：话术生成}};
        \draw[msg] (12.0,-6.1) -- node[above,font=\scriptsize]{调用remind\_delivery} (16.0,-6.1);
        \draw[msg] (16.0,-6.8) -- node[above,font=\scriptsize]{读session上下文+人设+用户画像} (8.0,-6.8);
        \draw[retmsg] (8.0,-7.3) -- node[above,font=\scriptsize]{chat\_session+user\_profile} (16.0,-7.3);
        \draw[msg] (16.0,-8.0) -- node[above,font=\scriptsize]{LLM生成话术（Ollama）} (8.0,-8.0);
        \draw[retmsg] (8.0,-8.5) -- node[above,font=\scriptsize]{个性化话术文本} (16.0,-8.5);
        \draw[retmsg] (16.0,-9.0) -- node[above,font=\scriptsize]{display\_line} (12.0,-9.0);

        % Phase 4: Push
        \node[note] at (13.8,-9.7) {\textit{阶段④：WebSocket推送}};
        \draw[msg] (12.0,-10.2) -- node[above,font=\scriptsize]{broadcast\_remind\_trigger\_to\_user} (16.0,-10.2);
        \draw[retmsg] (16.0,-10.8) -- node[above,font=\scriptsize]{\texttt{type:remind\_trigger}（含生成话术）} (20.0,-10.8);
        \draw[retmsg] (16.0,-11.4) -- node[above,font=\scriptsize]{\texttt{chunk\_audio}（可选）+WAV（MiMo语音）} (20.0,-11.4);
        \node[note,anchor=west] at (20.1,-11.9) {\scriptsize 米色气泡展示；可选语音朗读};

        % Offline user note
        \node[note] at (0.2,-12.8) {\textit{离线用户：}记录保持is\_triggered=0，重连时flush\_pending\_reminders补发};

        \foreach \x in {0,4.0,8.0,12.0,16.0,20.0}
            \fill (\x,-14.0) circle (2.5pt);
    \end{tikzpicture}
    }
    \caption{主动关怀系统端到端时序图（抽取→入库→调度→话术生成→WebSocket推送；离线用户重连后补发）}
    \label{fig:care_seq}
\end{figure}

\section{实现边界说明}

\textbf{已完整实现：}\texttt{remind\_trigger}库表与REST接口、常驻扫描调度器（\texttt{remind\_trigger\_scheduler.py}）、WebSocket推送帧\texttt{type:remind\_trigger}、\texttt{remind\_delivery}触发时LLM话术生成（含人设+用户画像+关联对话上下文拼装）、对话抽取入库（\texttt{remind\_extract.py}，绑定\texttt{session\_id}）、\texttt{scan-now}诊断接口与统计字段、MiMo可选语音朗读、离线用户重连补发（\texttt{flush\_pending\_reminders\_for\_connection}）。

\textbf{可继续扩展：}生成稿落库审计（当前生成稿不写回数据库）；更细粒度闹钟与窗口策略（当前为轮询扫描）；知识图谱式记忆召回（当前以MySQL+Redis结构化为主）；多连接策略统一（当前推送至该user\_id的所有在线连接）。

% ====================================================
% 第七章 结论
% ====================================================
\chapter{结论}

\section{工作总结}

本文面向高校学生日常低门槛情感交互练习与原型验证场景，交付浏览器端Live2D数字人系统一套。实现规模可概括为：9张MySQL核心表、61条REST路由；前端基于Vite组织的模块化源码与Cubism运行时；后端FastAPI编排Ollama双模型；朗读以小米MiMo为当前默认实现，GPT-SoVITS保留于仓库；完成ZIP模型入库与分层记忆工程化，含Redis列表与字符串、MySQL周期概要以及后台定时写入周期概要的任务；主动关怀系统实现从对话抽取到WebSocket推送的完整闭环，含后台常驻调度器、LLM个性化话术生成、MiMo可选语音朗读与离线补发机制。

\section{创新点}

\begin{enumerate}
    \item \textbf{全链路情感交互编排：}在单仓库内打通``语音识别—文本生成—动作决策—语音合成—角色呈现—会话沉淀''等可运行步骤。语音朗读以\texttt{chunk\_audio}元数据搭配二进制WAV格式，统一并入\textbf{对话WebSocket}（\texttt{/ws/chat}）通道，前端通过保序入队机制结合音频能量驱动数字人口型同步，有效解决``字幕与语音脱节''问题。
    \item \textbf{双模型分工的可控交互接口：}采用Ollama双模型解耦设计，将``自然语言台词生成''与``表情/动作结构化输出''分离为两次独立调用，由服务端完成格式校验与异常回退。相较单模型混合生成方案，该模式显著降低格式漂移风险，支持独立调参与日志溯源。
    \item \textbf{编排层主导的三级分层记忆：}构建瞬时（Redis List，最近$K$轮）、短期（过期轮次规则压缩）、长期（MySQL周期概要+Redis缓存）三级记忆机制，由服务端按固定路径读写，通过后台定时任务批量更新长期摘要。相较MemGPT``模型自主调度外部存储''思路，本文采用确定性编排机制，在本科研究周期内保障数据路径可追溯、可配置、可审计。
    \item \textbf{时间驱动的主动关怀完整闭环：}实现从对话内信息抽取（\texttt{remind\_extract}）、触发器持久化、常驻调度器、LLM个性化话术生成（\texttt{remind\_delivery}，拼装人设+用户画像+关联对话上下文）到WebSocket推送的全链路闭环，使系统具备超越被动应答的时间感知与主动情感关怀能力。离线用户重连后自动补发机制进一步保证关怀消息的可达性。
    \item \textbf{私有资源与会话数据一体化底层支撑：}基于用户登录鉴权，通过MySQL八表结构完成用户、数字人人设、会话、长期记忆、触发器等实体的标准化建模；Live2D资源包依托MinIO分布式存储，通过预签名URL实现浏览器端私有安全加载，配套ZIP批量入库与包级索引删除等一致性处理方案。
\end{enumerate}

\section{研究不足与未来工作}

\textbf{研究不足：}本文未开展系统化并发压力测试与大规模用户验证实验；语音识别模块依赖DashScope服务，存在第三方供应商绑定问题；合成时延相关性能数据仅基于开发环境少量观测，未留存原始实验数据，无法构成严格规范的对照实验；主动关怀的生成话术当前不落库审计，难以系统性评估话术质量。

\textbf{短期工作}优先围绕体验与稳定性：将关怀生成稿落库以支持质量审计；补充时间信息抽取规则与误触发抑制；完善鉴权与HTTPS部署清单；在用户授权前提下接入课程表、校历考试信息与轻量日记输入。\textbf{长期方向}包括探索检索增强RAG与知识图谱式记忆召回；在明确同意机制下评估社交动态等外部信息接入价值；逐步引入人格风格控制与动态养成策略，并用主观MOS与用户体验量表进行评估。

\textbf{深远期探索方向}着眼于使数字角色从``会对话的程序''向``有生命感的虚拟个体''演进，初步设想以下三个维度：

\textbf{其一，长期记忆的自然衰退机制。}人类记忆并非无损存档，而是随时间推移与重要性权重动态衰减；将这一机制引入分层记忆系统，可在长期概要层为每条记录附加时间戳与重要度评分，并由后台任务周期性对低权重、时间久远的条目执行软遗忘或压缩降维，使角色对久远细节的``记忆''自然模糊，从而获得更接近真实关系的连续性体验。

\textbf{其二，AI 辅助生成角色专属表情与动作资源。}现阶段系统依赖人工预制的表情与动作文件，能够表达的情绪粒度受资产数量限制。未来可深入研究 Live2D Cubism Editor 的网格变形原理与参数曲线结构，探索以生成模型自动合成符合目标角色骨骼与绑定约束的新表情（\texttt{.exp3.json}）和动作片段（\texttt{.motion3.json}），以此突破人工建模的产能瓶颈，为长期角色个性化运营提供可持续的资产来源。

\textbf{其三，角色自主情节安排与 AI 生成瞬间动画。}若赋予角色``自己的生活''——例如在非对话时段根据预设世界观自主规划当日行程、触发情节事件并将这些事件以简短独白或主动推送的形式呈现——将使用户感受到陪伴的真实性，而不仅仅是``被动等待回复的助手''。在此基础上，进一步结合扩散模型或轻量级 2D 动画生成技术，为关键情节自动合成带叙事感的短帧动画，有望在保持 Live2D 轻量部署优势的前提下，为用户提供更具代入感的情感共鸣体验。上述方向均超出本科阶段工程边界，留待后续研究深化。

\clearpage
\begin{thebibliography}{99}
\bibitem{MOE2023} 教育部等十七部门. 全面加强和改进新时代学生心理健康工作专项行动计划2023--2025年[Z], 2023.
\bibitem{NMHReport2023} 傅小兰, 张侃, 陈雪峰, 等. 中国国民心理健康发展报告2021--2022[R]. 北京: 社会科学文献出版社, 2023.
\bibitem{WHO2022} World Health Organization. World Mental Health Report: Transforming Mental Health for All[R]. Geneva: WHO, 2022.
\bibitem{Picard1997} Picard R W. Affective Computing[M]. Cambridge, MA: MIT Press, 1997.
\bibitem{Bickmore2005} Bickmore T W, Picard R W. Establishing and maintaining long-term human-computer relationships[J]. ACM Transactions on Computer-Human Interaction, 2005, 12(2): 293-327.
\bibitem{CharacterAIAbout} Character.AI. About Character.AI[EB/OL]. \url{https://character.ai/about/}.
\bibitem{ReplikaOfficial} Replika. Replika Official Website[EB/OL]. \url{https://replika.com/}.
\bibitem{live2d} Live2D Inc. Cubism SDK Manual -- Web[EB/OL]. \url{https://docs.live2d.com/cubism-sdk-manual/top/}.
\bibitem{Atkinson1968} Atkinson R C, Shiffrin R M. Human memory: A proposed system and its control processes[M]//Spence K W, Spence J T. The Psychology of Learning and Motivation. New York: Academic Press, 1968: 89-195.
\bibitem{MemGPT2023} Packer C, Fang V, Patil S, et al. MemGPT: Towards LLMs as operating systems[EB/OL]. arXiv:2310.08560, 2023. \url{https://arxiv.org/abs/2310.08560}.
\bibitem{Fitzpatrick2017} Fitzpatrick K K, Darcy A, Vierhile M. Delivering cognitive behavior therapy to young adults with symptoms of depression and anxiety using a fully automated conversational agent (Woebot): A randomized controlled trial[J]. JMIR Mental Health, 2017, 4(2): e19.
\bibitem{ShumXiaoice2020} Zhou L, Gao J, Li D, Shum H-Y. The design and implementation of XiaoIce, an empathetic social chatbot[J]. Computational Linguistics, 2020, 46(1): 53-93.
\bibitem{Roller2021} Roller S, Dinan E, Goyal N, et al. Recipes for building an open-domain chatbot[J]. arXiv preprint arXiv:2004.13637, 2021.
\bibitem{Zhao2023Survey} Zhao W X, Zhou K, Li J, et al. A survey of large language models[EB/OL]. arXiv:2303.18223, 2023. \url{https://arxiv.org/abs/2303.18223}.
\bibitem{AliyunAvatarProduct} 阿里云. 虚拟数字人产品介绍[EB/OL]. \url{https://www.aliyun.com/product/ai/avatar}.
\bibitem{SenseAvatarProduct} 商汤科技. 日日新·如影 SenseAvatar 产品页面[EB/OL]. \url{https://www.sensetime.com/cn/product-detail?categoryId=51134279}.
\bibitem{NeteaseFuxiAgent} 网易伏羲. 虚拟智能体产品页[EB/OL]. \url{https://fuxi.163.com/productDetail/agent}.
\bibitem{Li2021EmotionDialogueCN} 李明, 王昊奋, 刘挺. 开放域情感对话系统研究进展[J]. 中文信息学报, 2021, 35(10): 1-15.
\bibitem{Chen2022EmpatheticDialogueCN} 陈思宇, 黄萱菁, 林鸿飞. 共情对话生成技术综述[J]. 计算机科学, 2022, 49(6): 30-42.
\bibitem{Sun2022DialogueMemoryCN} 孙悦, 何婷婷, 郑伟. 面向多轮对话的记忆增强方法研究综述[J]. 计算机工程, 2022, 48(11): 1-12.
\bibitem{He2023UserProfileCN} 何佳, 陈琳, 蒋勇. 用户画像驱动的个性化对话服务研究[J]. 情报理论与实践, 2023, 46(8): 97-104.
\bibitem{ReAct2023} Yao S, Zhao J, Yu D, et al. ReAct: Synergizing Reasoning and Acting in Language Models[C/OL]//ICLR 2023. \url{https://arxiv.org/abs/2210.03629}.
\bibitem{GenerativeAgents2023} Park J S, O'Brien J C, Cai C J, et al. Generative Agents: Interactive Simulacra of Human Behavior[C]//UIST 2023. New York: ACM, 2023.
\bibitem{Lewis2020} Lewis P, Perez E, Piktus A, et al. Retrieval-Augmented Generation for Knowledge-Intensive NLP Tasks[C]//NeurIPS. 2020, 33: 9459-9474.
\bibitem{OpenCLaW2024} OpenCLaW Contributors. OpenCLaW: Open-source Framework for Collaborative Language Agent Workflows[EB/OL]. 2024. \url{https://github.com/openclaw/openclaw}.
\bibitem{Harness2024} EleutherAI. Language Model Evaluation Harness[EB/OL]. 2024. \url{https://github.com/EleutherAI/lm-evaluation-harness}.
\bibitem{rfc6455} Fette I, Melnikov A. The WebSocket Protocol: RFC 6455[S/OL]. IETF, 2011. \url{https://www.rfc-editor.org/rfc/rfc6455}.
\bibitem{mdn-ws} MDN Web Docs. WebSocket API[EB/OL]. \url{https://developer.mozilla.org/zh-CN/docs/Web/API/WebSocket}.
\bibitem{ollama} Ollama Documentation[EB/OL]. \url{https://github.com/ollama/ollama}.
\bibitem{gptsovits} RVC-Boss. GPT-SoVITS[EB/OL]. \url{https://github.com/RVC-Boss/GPT-SoVITS}.
\bibitem{mimoV25Docs} 小米 MiMo. Speech synthesis v2.5 usage guide[EB/OL]. Xiaomi MiMo Platform. \url{https://platform.xiaomimimo.com/docs/zh-CN/usage-guide/speech-synthesis-v2.5}.
\bibitem{mysql80} Oracle. MySQL 8.0 Reference Manual[EB/OL]. \url{https://dev.mysql.com/doc/refman/8.0/en/}.
\bibitem{redis} Redis Ltd. Redis Documentation[EB/OL]. \url{https://redis.io/docs/}.
\bibitem{minio} MinIO, Inc. MinIO Object Storage Documentation[EB/OL]. \url{https://min.io/docs/minio/}.
\bibitem{mdn-audio} MDN Web Docs. Web Audio API 指南[EB/OL]. \url{https://developer.mozilla.org/zh-CN/docs/Web/API/Web_Audio_API}.
\bibitem{vite} Vite Team. Vite Next Generation Frontend Tooling[EB/OL]. \url{https://vitejs.dev/}.
\bibitem{fastapi} FastAPI Documentation[EB/OL]. \url{https://fastapi.tiangolo.com/}.
\bibitem{ffmpeg} FFmpeg Team. FFmpeg Documentation[EB/OL]. \url{https://ffmpeg.org/documentation.html}.
\bibitem{CAICTDigitalHuman2023} 中国信息通信研究院. 虚拟数字人发展白皮书2023年[R]. 北京: 中国信通院, 2023.
\bibitem{Wang2023CampusDigitalHumanCN} 王磊, 刘洋, 赵倩. 高校场景下虚拟数字人应用模式与实现路径研究[J]. 现代教育技术, 2023, 33(9): 112-120.
\bibitem{Liu2021StreamingSystemCN} 刘志强, 周明, 陈震. 面向实时交互的流式对话系统架构设计[J]. 软件导刊, 2021, 20(12): 15-22.
\bibitem{Wu2024AIEducationCN} 吴岩, 李志民. 人工智能赋能高等教育数字化转型的关键问题与治理框架[J]. 中国高教研究, 2024(3): 12-19.
\bibitem{CharacterAIProduct} Character.AI. Character.AI Product Site[EB/OL]. \url{https://character.ai/}.
\bibitem{ReplikaHelp} Replika. What is Replika?[EB/OL]. \url{https://help.replika.com/hc/en-us/articles/115001070951-What-is-Replika}.
\bibitem{HololiveProduction} COVER Corp. hololive production[EB/OL]. \url{https://hololive.hololivepro.com/en/}.
\bibitem{Toolformer2023} Schick T, Dwivedi-Yu J, Dessi R, et al. Toolformer: Language Models Can Teach Themselves to Use Tools[C]//NeurIPS, 2023.
\bibitem{Rashkin2018} Rashkin H, Smith E M, Li M, Boureau Y L. Towards Empathetic Open-domain Conversation Models[C]//ACL. Melbourne: ACL, 2018: 537-549.
\end{thebibliography}

\clearpage
\chapter*{致谢}
\addcontentsline{toc}{chapter}{致谢}

大学四年接近尾声，本次毕业设计既是对所学人工智能与软件开发知识的综合运用，也是我第一次把``能跑的工程''与``能写得清的论文''对齐起来的完整演练。

回顾整个开发过程，令我印象最深的是两个系统的落地经历：分层记忆的工程化，让我第一次真切地感受到``把历史信息从Redis搬进MySQL、再在合适的时机取出注入模型''这件事在工程上并不简单，需要反复在``数据在哪里、什么时候读、读多少''之间做取舍；主动关怀系统的闭环，则让我理解了``定时触发''在工程上并不是精确的闹钟，而是``到期+扫描发现''的两阶段语义——这种对工程细节的理解，是单纯学理论学不到的。

语音链路曾依赖本地GPT-SoVITS，实测合成耗时偏高，难以与对话首包体验兼顾；后以小米MiMo作为当前主线，又面临限流与参考音频重复上行的现实约束——最终将句末标点攒批阈值由1调整为4，以合并更长的文本段再发起单次合成，从而减少请求频次。这一过程让我体会到：论文中的``最优参数''未必总能直接落到线上环境，\textbf{部署约束与观测数据}同样构成设计的一部分。

首先要衷心感谢导师\textbf{汪柳萍}老师，在选题边界、工程可实现性与论文叙事结构上给予了耐心指导，特别是在如何把系统的差异化写得具体、如何界定``已完成''与``可扩展''方面，帮助我反复校正叙述重心。

感谢朋友们在联调、界面体验与文档校对上的帮助；感谢家人在日常学习与生活上的支持。最后感谢开源社区——Live2D、FastAPI、Ollama、Redis、MySQL、MinIO以及本项目所依赖的其他库——没有这些可复用的基础设施，本科阶段难以在有限时间内完成端到端系统。

\end{document}